% =====================================================================
% Введение.
% Строгий порядок по правкам друга-ревьюера
%   1. Актуальность работы
%   2. Новизна
%   3. Предмет исследования
%   4. Объект исследования
%   5. Цель работы (без слова «оценка», только «количественная оценка»)
%   6. Аналитический обзор предлагаемых решений (краткий)
%   7. Практическая значимость
%   8. Теоретическая значимость
%   9. Методы исследования
%  10. Задачи работы
%  11. Структура работы (с числом таблиц / рисунков / страниц / формул)
%
% Объем 3--4 страницы.
% =====================================================================

\section*{Введение}
\addcontentsline{toc}{section}{Введение}

\textbf{Актуальность работы.} В 2023--2026 годах применение больших
языковых моделей (от англ. Large Language Models, LLM) перешло от
одиночных диалоговых сценариев к кооперативной работе нескольких
автономных агентов, образующих мульти-агентную систему
(от англ. Multi-Agent System, MAS). Открытые фреймворки (AutoGen,
ChatDev, MetaGPT, Magentic-One) показывают, что разделение труда
между специализированными агентами повышает качество решения
сложных задач, но каждый из них жестко фиксирует структуру
коммуникации на этапе проектирования; систематических данных об
оптимальной топологии для каждого класса задач в открытой литературе
нет. Параллельный интерес к интеграции человека в контур управления
ограничивается единичными паттернами без сравнения эффективности.
Сочетание этих обстоятельств определяет научную и прикладную
актуальность работы.

\textbf{Научная новизна.} В настоящей работе впервые проведено
прямое сравнение пяти базовых коммуникационных топологий
(звезда, цепочка, полносвязная, дебатная, иерархическая) на
едином программном стенде, на одной языковой модели и на одной
выборке задач из четырех различных классов. Предложена и
реализована мета-топология, переключающая активную базовую
конфигурацию в зависимости от текущей фазы решения и наблюдаемых
сигналов агентов. Введена таксономия из пяти специфичных для
адаптивных мульти-агентных систем метрик ---
$N_{\text{switch}}$, $r_{\text{guard}}$, $\gamma$,
$r_{\text{hurt}}$, $\Delta_{\text{oracle}}$, --- позволяющая
количественно характеризовать поведение мета-топологии. Впервые
сформулирована и экспериментально проверена матрица совместимости
пяти ролей человека в контуре управления с пятью базовыми
топологиями.

\textbf{Предмет исследования.} Коммуникационная топология
мульти-агентной системы больших языковых моделей, механизм
ее динамической адаптации в процессе решения задачи, а также
позиция человека в графе обмена сообщениями между агентами.

\textbf{Объект исследования.} Мульти-агентные системы на основе
больших языковых моделей, решающие сложные задачи различных
классов в кооперации со специализированными агентами и
участником-человеком.

\textbf{Цель работы.} Количественная оценка влияния выбора
коммуникационной топологии мульти-агентной системы, механизма
ее адаптации и позиции человека в контуре управления на
эффективность решения задач разных классов с одновременным
учетом качества решения, его стоимости и времени выполнения.

\textbf{Аналитический обзор предлагаемых решений.} В 2022--2026
годах выделяются три направления развития MAS LLM. Архитектурные
фреймворки (AutoGen, ChatDev, MetaGPT) фиксируют топологию и не
сравнивают альтернативы. Автоматическое проектирование графов
(G-Designer, GPTSwarm) подбирает структуру однократно и не
пересматривает ее. Динамическая адаптация (DyLAN, AgentVerse)
меняет состав команды, но не структуру связей. Вопрос роли человека
освещен фрагментарно, без количественной оценки ролей. Подробный
анализ --- в главе~\ref{sec:review}, сравнительная
таблица~\ref{tab:rev:methods} восемнадцати методов по шести
критериям.

\textbf{Практическая значимость.} Работа дает разработчикам MAS
количественные данные о том, какая топология эффективнее для
каждого из четырех рассмотренных классов задач и когда имеет
смысл адаптивное переключение. Программная инфраструктура
реализована как открытый фреймворк адаптивных топологий
(от англ. Adaptive Topologies for Multi-agent systems, ATM),
исходный код --- в репозитории GitHub (приложение~\ref{app:github}),
допускает повторное использование и расширение. Сформулированные
правила выбора топологии и роли человека под класс задачи применимы
при проектировании интеллектуальных ассистентов в программировании,
аналитике и принятии решений.

\textbf{Теоретическая значимость.} Введенная таксономия
адаптивно-специфичных метрик ($N_{\text{switch}}$, $r_{\text{guard}}$,
$\gamma$, $r_{\text{hurt}}$, $\Delta_{\text{oracle}}$) позволяет
количественно характеризовать поведение динамически перестраиваемых
MAS по единой шкале. Формализация MAS как кортежа из шести
компонентов с фазово-зависимой коммуникационной структурой
расширяет теоретическую базу; эмпирическая проверка гипотезы о
фазовой специфике топологий дает обоснование для дальнейших
теоретических работ.

\textbf{Методы исследования.} Комплекс взаимодополняющих методов:
аналитический обзор литературы по ключевым словам в открытых базах;
формализация MAS средствами теории графов и конечных автоматов;
программная реализация на Python с LangGraph; контролируемый
эксперимент с предварительной регистрацией гипотез и порогов
улучшения; статистическая обработка (дисперсионный анализ от англ.
Analysis of Variance, ANOVA; парные сравнения с поправкой
Бенджамини--Хохберга на уровне ложных открытий от англ. False
Discovery Rate, FDR; коэффициент Коэна; непараметрический
бутстрап с доверительными интервалами от англ. Confidence
Interval, CI); оракульный референс per-task best-of
для верхней границы качества адаптации; имитационное моделирование
участника-человека языковой моделью с подсказкой по роли.

\textbf{Задачи работы.} Для достижения сформулированной выше
цели в работе решаются следующие задачи.
\begin{enumerate}
  \item Провести аналитический обзор работ 2022--2026 годов по
        мульти-агентным системам больших языковых моделей,
        коммуникационным топологиям и интеграции человека в
        контур управления, выделить исследовательскую лакуну.
  \item Сформулировать формальную модель мульти-агентной системы
        как кортежа из множества агентов, множества ребер
        коммуникации, множества ролей агентов и человека,
        задачи и множества фаз решения.
  \item Разработать и реализовать программный фреймворк, в
        котором пять базовых топологий и адаптивная мета-топология
        реализованы в единой архитектуре на основе библиотеки
        LangGraph.
  \item Спроектировать систему оценочных метрик, включающую
        базовые показатели качества, стоимости и времени, а также
        специфические для адаптивной мета-топологии характеристики
        переключений, заблокированных решений маршрутизатора,
        стоимостного вклада маршрутизатора, доли регрессий
        качества и разрыва до верхнего потолка адаптации.
  \item Провести воронку из четырех экспериментов и
        подтверждающего прогона на четырех классах задач ---
        программирование, математические рассуждения, творческая
        генерация, анализ данных --- с предварительно
        зафиксированными гипотезами и порогами улучшения.
  \item Выполнить статистический анализ результатов и сформулировать
        ответы на четыре исследовательских вопроса (от англ.
        Research Question, RQ), описанных в
        разделе~\ref{sec:meth:funnel}.
  \item Выработать практические рекомендации по выбору топологии
        и роли человека в зависимости от класса задачи.
\end{enumerate}

\textbf{Структура работы.} Работа состоит из введения, пяти глав
основного содержания, раздела авторского вклада, заключения,
библиографического списка и четырех приложений (А, Б, В, Г).
Глава~1 --- аналитический обзор литературы и исследовательская
лакуна. Глава~2 --- формальная модель, шесть топологий, пять
ролей человека, метрики, корпус, дизайн воронки. Глава~3 ---
архитектура программной системы. Глава~4 --- результаты четырех
основных экспериментов и подтверждающих прогонов на дополнительном
семействе моделей. Глава~5 --- анализ, ответы на четыре RQ и
угрозы валидности. Раздел авторского вклада фиксирует оригинальные
элементы; заключение подводит итоги и формулирует направления
дальнейших исследований. Приложения --- функциональная схема,
глоссарий, состав репозитория, ссылка на исходный код.

В работе пять глав основного содержания, четыре приложения,
19~таблиц, 8~рисунков и~5~формул; общий объем работы с~введением,
заключением, библиографическим списком и~приложениями составляет
62~страницы; библиографический список --- 44~источника.
